% ======================================================================
%  50_appendix_compare.tex —— 附录 A：与三份强相关文献的结构化对照
%  只允许 \input、环境、preamble 已定义的宏；不加载宏包、不新定义命令。
% ======================================================================

\section{与三份强相关文献的结构化对照}
\label{sec:appendix-compare}

本附录是全书唯一的横向比较。它不问「哪一份文献更完整、更严格」，只问一个结构问题：底本那条 S1--S8 的八步主链（编号见骨架层），在三份强相关文献里分别被\textbf{保留、压缩、改写、替换或补全}成了什么？因此下面每一小节都按同一组栏目组织：进入点、被保留的步骤、被省略的步骤、新增的视角、未补齐的缺口。\pC 这样编排的代价是不按文献自己的章节顺序复述；好处是读者可以把三份文献放回同一条推理链上，看清它们各自占住哪一段。

四份文本的基本事实先立此存照（表 A.1），因为其中至少三处标题会误导读者。

\begin{center}
\small\textbf{表 A.1}\quad 本附录对照的四份文本\\[3pt]
\footnotesize
\begin{tabular}{@{}p{0.15\textwidth}p{0.34\textwidth}p{0.43\textwidth}@{}}
\toprule
\tabhead{文本} & \tabhead{版本与规模} & \tabhead{在本对照中的角色} \\
\midrule
KW06（底本） & arXiv:hep-th/0604151，2006 年 4 月；\emph{CNTP} 1 (2007)；231 页、12 节、146 条参考文献 \anchor{src p.1--232} & S1--S8 主链与 M1--M5 五个 miracle 的出处 \\
Witten 2008 & \emph{Mirror Symmetry, Hitchin's Equations, and Langlands Duality}；arXiv:0802.0999，2008 年 2 月；15 页、6 节 \anchor{0802.0999 p.1} & 把 S3--S4 改写为 Hitchin 纤维化上的镜像对称 \\
Witten 2009 & \emph{Geometric Langlands From Six Dimensions}；arXiv:0905.2720，2009 年；30 页、5 节；2008 年 6 月 Raoul Bott 纪念会讲稿 \anchor{0905.2720 p.1} & 把 S1、S7 压成陈述，并换成六维起点 \\
Frenkel 2003 & \emph{Lectures on the Langlands Program and Conformal Field Theory}；arXiv:hep-th/0512172，2005 年 12 月；128 页、9 节；2003 年 3 月 Les Houches 讲稿 \anchor{Frenkel p.1} & 从顶点算子代数与共形场论进入同一对象 \\
Setter 2013 & \emph{Topological Quantum Field Theory and the Geometric Langlands Correspondence}；Caltech 博士论文，2013 年（2012-08-24 答辩）；169 页；导师 A. Kapustin \anchor{Setter p.i} & 把 K--W 的框架往下推广到 3d 与面算子范畴 \\
\bottomrule
\end{tabular}
\end{center}

\begin{warn}
\textbf{三个篇目陷阱。}（1）有些编目把「Geometric Langlands From Six Dimensions」记作 0802.0999；实际 0802.0999 是 15 页 6 节的 \emph{Mirror Symmetry, Hitchin's Equations, and Langlands Duality}（\anchor{0802.0999 p.1}），六维讲稿是 30 页 5 节的 0905.2720（\anchor{0905.2720 p.1}）。本附录对二者分别设锚点，绝不把六维内容记在 0802.0999 名下。\pA（2）0905.2720 里的 ``brane'' 多数指 M5 膜、D4 膜、D6 膜，而不是 2006 年的 A/B-膜；同一个词在两篇里不是同一类对象。\pA（3）Setter 的论文标题含 geometric Langlands，但正文既不证明 GL，也不定义 Hecke 作用（见 \S\ref{sec:compare-setter}）。\pA
\end{warn}

阅读约定：\pA 表示原文或对照文献明写；\pB 表示原文只是暗示、由我们重建；\pC 表示我们的判断。锚点 \anchor{...} 一律给印刷页（PDF 页与印刷页的换算见底本说明）。短引一律不超过 200 字符，只用于无法转述的关键句；表 A.2--A.5 每格除标 \pB、\pC 者外均为 \pA。

% ----------------------------------------------------------------------
\subsection{与 Witten 2008/2009 对照：八步主链剩下几句话}
\label{sec:compare-witten}

\subsubsection*{进入点：从「推导 GL」改为「GL 是 4d 对偶的范畴化陈述」}

2009 年那篇讲稿的做法，可以用它自己的两句话概括。第一句说明扭化与 GL 的关系：\pA
\begin{origquote}
The equivalence between this third twisting for the two groups $G$ and $G^\vee$ \ldots\ leads precisely at the level of ``categories,'' that is for two-manifolds, to the geometric Langlands correspondence. \anchor{0905.2720 p.4}
\end{origquote}
第二句说明 2009 的公理层级：一个 $d$ 维量子场论给 $d$ 维流形一个配分函数、给 $d-1$ 维流形一个 Hilbert 空间、给 \textbf{$d-2$ 维流形一个范畴}（\anchor{0905.2720 p.1}）。\pA 于是 2006 用八步造出来的东西，在 2009 里成了「$d=4$ 时范畴落在二维，而第三种扭化在 $G$ 与 $\Ldual$ 之间给出范畴等价」这一句。这是本附录要解释的第一个现象：不是 2009 把 2006 的证明做短了，而是它\textbf{不再需要}那条证明——它只需要结论，把证明换成公理。

2008 年那篇 15 页短文走的是另一条路。它不引入六维，而是把 2006 的 S3--S4 重新讲成 Hitchin 纤维化上的镜像对称（\anchor{0802.0999 p.2}），并在最后明确说四维规范论是它的输入：\pA
\begin{origquote}
Taking this as input \ldots\ one arrives at the twisted version of $N=4$ super Yang-Mills theory that was the starting point in Kapustin \& Witten (2007). \anchor{0802.0999 p.14}
\end{origquote}
这句话把三篇的分工写得很清楚：2008 从 2006 的扭化理论出发、往几何方向走；2009 从六维出发、往四维走；两者都把 2006 的 4d 推导当作既成事实。

\subsubsection*{映射表：2006 主张 $\to$ 2008/2009 表述}

\begin{center}
\small\textbf{表 A.2}\quad S1--S8 与 M1--M5 在两篇 Witten 中的命运\\[3pt]
\footnotesize
\begin{tabular}{@{}p{0.11\textwidth}p{0.25\textwidth}p{0.31\textwidth}p{0.13\textwidth}@{}}
\toprule
\tabhead{2006 主张} & \tabhead{2006 的表述（锚点）} & \tabhead{2008/2009 的对应（锚点）} & \tabhead{处理} \\
\midrule
S1 对偶群 & $G$ 与 $\Ldual$ 的规范论在 $\tau\mapsto-1/(n_g\tau)$ 下等价；Table 1 列对偶群 \anchor{src p.3}，\S2 \anchor{src p.6} & 0905 \S1.2 一句 GNO 对偶群 \anchor{0905.2720 p.3}；\S5.5.3 只处理 simply-laced \anchor{0905.2720 p.27}；0802 \S2 仍用 GNO/Langlands 对偶 \anchor{0802.0999 p.2} & 压缩，并加限制 \\
S2 扭化族 & $N=4$ 有一族由 $\PP^1$ 参数化的拓扑扭化，整族保持拓扑性（M1）\anchor{src p.16}、\anchor{src p.28} & 0905 p.3--4：三种扭化，第三种与 GL 相关；扭化参数破坏 $G\leftrightarrow\Ldual$ 对称 \anchor{0905.2720 p.3} & 压缩为两段 \\
S3 紧化到 $C$ & 4d TFT 在黎曼面 $C$ 上紧化，得靶为 $\MH$ 的 2d 拓扑 sigma 模型 \S4 \anchor{src p.38} & 0905 p.1--2 只在公理层说「给 2-流形一个范畴」\anchor{0905.2720 p.1}；显式紧化改为六维：$M_6=M_4\times T^2$ \anchor{0905.2720 p.8}、圆丛 \anchor{0905.2720 p.15} & 省略并改向 \\
S4 2d TFT 与镜像 & 4d $S$-对偶在 2d 变成镜像对称；$\MH(G)$ 与 $\MH(\Ldual)$ 在复结构 $J$ 下互为镜像（M2）\S5.4 \anchor{src p.71}、\S5.5 \anchor{src p.74} & 0905 不提 B/A-模型、不提镜像对称（``B-model'' 0 次）；只保留 $\tau\mapsto-1/\tau$ \anchor{0905.2720 p.4}、\anchor{0905.2720 p.10} & 0905 省略；0802 \S2--3 补回 \\
S5 线算子$=$膜自函子 & 尊重拓扑对称的线算子 $X$ 在膜上有自然作用，给出膜范畴的自函子 \S6 \anchor{src p.78}、\S6.6 \anchor{src p.100} & 0905 \S2 改写为「$p$-流形算子」：Wilson 与 't Hooft 都是 1-流形算子 \anchor{0905.2720 p.5} & 改写为缺陷维数 \\
S6 特征膜$=$局部系统 & Wilson 算子的本征膜 $\leftrightarrow$ $\pi_1(C)\to\LdualC$ 的表示 \S7--8 \anchor{src p.103}、\anchor{src p.116} & 0905 无 eigenbrane；只以「$Y_{\Ldual}(C)$ 上凝聚层 $\leftrightarrow$ $Z_G(C)$ 上 $D$-模」一句代替 \anchor{0905.2720 p.3} & 省略 \\
S7 't Hooft $=$ Hecke & 't Hooft 算子对应 Hecke 算子，经 Bogomolny 方程与单极子冒泡证明 \S9--10 \anchor{src p.123}、\anchor{src p.148} & 0905 p.6 一句：the 't Hooft operators correspond to the usual geometric Hecke operators \anchor{0905.2720 p.6}；无 Bogomolny、无冒泡 & 压缩为一句 \\
S8 c.c.\ 膜$\to$\mbox{$D$-模} & $\MH$ 上的典范余各向同性 A-膜使 magnetic eigenbrane 自然对应 $D$-模 \S11 \anchor{src p.184}、\S11.1 \anchor{src p.186} & 0905 无 c.c.\ 膜；``D-module'' 仅在 p.3 出现两次 \anchor{0905.2720 p.3}；0802 \S3.1 引 [KW \S11] 补回 \anchor{0802.0999 p.5} & 0905 省略；0802 引用 \\
新增：六维 $(0,2)$ & 2006 全文无六维 $(0,2)$ 内容 & 0905 \S3--4：$H=\star H$ 的 gerbe 理论 \anchor{0905.2720 p.6}；$M_6=M_4\times T^2$ 给 4d 规范论、$\tau=\tau(T^2)$ \anchor{0905.2720 p.8}；非阿贝尔 type $G$ 理论 \anchor{0905.2720 p.9} & 新增 \\
新增：M5 膜与复曲面 & 无 & 0905 p.18：$U(n)$ 的 $(0,2)$ 即 $n$ 个平行 M5 膜，$n$ 个 D4 给 $N=4$ SYM \anchor{0905.2720 p.18}；\S5 复曲面版 GL，TN$_k$ 上 BPS 态 $=\bigotimes_j W_j$ \anchor{0905.2720 p.25} & 新增 \\
\bottomrule
\end{tabular}
\end{center}

\subsubsection*{压缩的算术：不是缩短，而是只留下两端}

表 A.2 的第四列只出现三种取值：压缩、省略、新增。八步之中，以\textbf{陈述}形式留在 2009 的只有 S1 与 S7（各一句），以\textbf{公理}形式留下的是 S3 与 S5 背后的范畴语言，其余全部消失：S2 的技术内容（$\PP^1$ 上的一族扭化、整族拓扑性的证明）、S4 的镜像对称、S6 的特征膜、S8 的 $D$-模构造。\pA 最刺眼的证据就是上引 0905 p.4 那句：一句话接管了 2006 用 §3--§11（约 170 页）建立的结论。\pA 用主链的语言说，2009 保留的是\textbf{两端}（S1 的对偶群、S7 的 't Hooft$=$Hecke），中间六步——扭化族、紧化、2d TFT、膜算子、特征膜、$D$-模——被折叠进「范畴」这个公理里。\pC 这并不是作者偷懒：2009 的目标读者是 Langlands 方向的数学家，讲稿的任务是给出六维框架，而不是重演 2006 的 231 页。\pA 但对照时必须说清楚：压缩的代价是\textbf{可重建性}——只看 0905 的读者无法知道 $\MH$ 的哪个复结构让理论变成 A-模型（2006 §4.2 的工作），也无法知道 Hecke 算子为什么等于 't Hooft 算子（2006 §9--10 的 Bogomolny 方程与单极子冒泡）。\pB

一个直接的文本证据是术语统计：0905 全文 ``Hitchin'' 只出现在参考文献（[10] 即 0802.0999，以及 Beilinson--Drinfeld 的书名）中，``B-model'' 出现 0 次，``Hecke'' 出现 5 次（p.6 一句正文，其余在文献表），``D-module'' 出现 2 次。\pA 因此 2008 年那篇 15 页短文不是 2009 的草稿，而是 2006 与 2009 之间\textbf{唯一仍以 Hitchin 纤维化为主角}的 Witten 文本；它自己也承认 4d 起点是输入（上引 p.14）。\pB 三篇合起来构成一条「长推导（2006）$\to$ 几何摘要（2008）$\to$ 六维纲领（2009）」的谱，而不是三次重复。\pC

\subsubsection*{新增视角的实质：换起点，不是补证明}

六维视角的价值不在于把 2006 的论证补严，而在于\textbf{替换起点}：4d 的电磁对偶在 2009 里不是假设，而是六维 $(0,2)$ 理论的推论。\pA 作者对此有一句非常坦率的话：\pA
\begin{origquote}
Its existence is probably our best explanation of electric-magnetic duality -- and therefore, in particular, of geometric Langlands duality. \anchor{0905.2720 p.9}
\end{origquote}
这里的 ``its existence'' 指 type $G$ 的六维 $(0,2)$ 理论。\pA 由此派生出三个新对象。（i）$M_6=M_4\times T^2$ 的模参数就是 4d 耦合：$\tau=\tau(T^2)$，$S$-对偶是 $T^2$ 的模对称性（\anchor{0905.2720 p.8}）；两种紧化次序给出 $I_4=\frac{S}{8\pi R}\int\Tr F\wedge\star F$ 与 $\frac{R}{S}$，于是 $\tau\mapsto-1/\tau$ 成为恒等式（\anchor{0905.2720 p.10}）。\pA（ii）M5 膜给出具体模型：$U(n)$ 的 $(0,2)$ 理论就是 $n$ 个平行 M5 膜，$n$ 个 D4 膜的低能理论是 $N=4$ 的 $U(n)$ SYM（\anchor{0905.2720 p.18}）。\pA（iii）Taub-NUT 上的 BPS 态空间等于 $k$ 个 level-1 WZW 共形块之张量积（\anchor{0905.2720 p.25}），这条等式把规范论与表示论（Nakajima、Braverman--Finkelberg、Licata）接上，是 2009 与数学文献之间最实在的一座桥。\pA

所谓「't Hooft 算子的六维起源」，严格说应改写为「\textbf{缺陷维数的重新定位}」：2009 把 Wilson 算子与 't Hooft 算子都定义为 1-流形算子（\anchor{0905.2720 p.5}），于是 $d-2=2$ 的 2-流形算子自然出现，而后者恰好给出分歧的 QFT 框架——「to get a class of two-manifold operators that is invariant under electric-magnetic duality, one must incorporate certain quantum parameters」\anchor{0905.2720 p.6}。\pA 这一句是 2006 与 2009 之间真正的接口：它把 2006 明确留在门外的分歧，接上了 2006 自己造出的工具（'t Hooft 算子与量子参数）。\pB

\subsubsection*{「绕开 4d 扭化」的真相}

2009 的叙事是「六维在先、四维在后」，但\textbf{逻辑上}它并没有绕开 2006 的扭化：p.3--4 仍以「第三种扭化的两个群之间的等价」为前提来陈述 GL（上引第一句），它真正重推的只是 $S$-对偶本身的六维起源。\pC 换句话说，「绕开」是修辞上的（把 4d 看作 6d 的推论），不是数学上的（GL 的范畴等价仍是 4d 扭化的结论）。\pC 这也解释了为什么 2009 对 2006 的 S2--S8 只能压缩而不能重写：那些步骤的技术内容（膜、Bogomolny 方程、$D$-模）在六维框架里没有对应物，作者也没有试图给它们六维版本。\pC 唯一的例外是 §5 的「复曲面版几何 Langlands」（\anchor{0905.2720 p.15}），但那是一个\textbf{新陈述}，不是对 2006 陈述的重新证明。\pA

\subsubsection*{没有补齐的五个缺口}

第一，仍只做 simply-laced：§5.5.3 明说此前假设 $G$ 与 $\Ldual$ 同李代数，非 simply-laced 要靠外自同构 $\lambda$-twist 并引 Vafa，未展开（\anchor{0905.2720 p.27}）。\pA 第二，仍不做证明：全文没有定理，§5 的 BPS 态计数是证据而不是对偶的证明（\anchor{0905.2720 p.25}）。\pA 第三，六维理论本身没有被构造：它被明确称为孤立非高斯共形场论、无经典 PDE 描述（\anchor{0905.2720 p.9}）；因此「用 6d 解释 4d 对偶」是把解释建立在一个未被构造的对象上。\pA（前半）\pC（后半） 第四，§5.5.4 中 level-$r$ WZW 与 $r$ 阶共轭类的对应，作者明说 ``this will not be attempted here''（\anchor{0905.2720 p.29}）。\pA 第五，2006 自陈的第二个 gap——不理解 2d CFT 对 GL 的用处（\anchor{src p.5}）——只被「六维共形块 $\approx$ level-1 WZW」的类比部分回应（\anchor{0905.2720 p.14}、\anchor{0905.2720 p.25}），并没有给出 Frenkel 式的 Hecke 特征层构造。\pC 所以 2009 补齐的是「对偶从哪来」，没有补齐「GL 如何证」与「分歧如何范畴化」。\pC

\paragraph{自检问题}
\begin{enumerate}
\item 如果只读 0905，你能独立重建 S5--S8 中的哪一步？提示：先看它是否出现过 ``brane category'' 与 ``$D$-module''。\pC
\item 2009 把 4d 对偶降为 6d 的推论，为什么这不等于「证明了」4d 对偶？\pC
\item 为什么 0802.0999 的 15 页是理解 2006 S3--S4 的捷径，却不能替代 2006 的 §9--§11？\pC
\end{enumerate}

% ----------------------------------------------------------------------
\subsection{与 Frenkel 2003 对照：两条入口，一套对象}
\label{sec:compare-frenkel}

\subsubsection*{进入点：一个解释为什么，一个构造是什么}

Frenkel 自己把 K--W 的贡献概括成三步：4d 规范论约化到 2d、't Hooft 算子变成 Hecke 算子、't Hooft 本征膜变成 Hecke 特征层。\pA 原文是：\pA
\begin{origquote}
a t'Hooft ``eigenbrane'' becomes after the reduction a Hecke ``eigensheaf'' \ldots\ Kapustin and Witten show that Langlands duality is closely related to S-duality of QFT. \anchor{Frenkel p.3}
\end{origquote}
这段转述值得注意：Frenkel 把 K--W 读成「约化 $+$ 算子对应 $+$ 膜与层的对应」，而不是读成「扭化 $+$ 紧化」。\pB 他自己的入口则写在 p.4：几何 Langlands 预言「带联络的全纯 $\Ldual$-丛 $\mapsto$ $\mathrm{Bun}_G$ 上 Hecke 特征层」，而 2d 共形场论正是构造 Hecke 特征层的机器（\anchor{Frenkel p.4}）。\pA 一句话概括：K--W 解释这条对应\textbf{为什么}存在，Frenkel 构造这条对应\textbf{是什么}。\pC

\subsubsection*{对照表：对象、工具、对偶、分歧}

\begin{center}
\small\textbf{表 A.3}\quad 两条路线的入口对照\\[3pt]
\footnotesize
\begin{tabular}{@{}p{0.13\textwidth}p{0.34\textwidth}p{0.35\textwidth}@{}}
\toprule
\tabhead{主题} & \tabhead{Frenkel 2003} & \tabhead{K--W 2006} \\
\midrule
对象 A 侧 & Hecke 特征层 $=$ $\mathrm{Bun}_G$ 上满足 Hecke 本征性的 $D$-模 \anchor{Frenkel p.4}、\anchor{Frenkel p.64} & magnetic eigenbrane，经 c.c.\ A-膜对应 $\MC$ 上 $D$-模（S8）\anchor{src p.184}、\anchor{src p.186} \\
对象 B 侧 & 带联络的全纯 $\Ldual$-丛，即局部系统 \anchor{Frenkel p.4} & $Y_{\Ldual}(C)$ 上的平坦丛与凝聚层（S6）\anchor{src p.116} \\
工具 & 临界 level $k=-h^\vee$ 的仿射代数、共形块、opers、Hitchin 系统量子化 \S7--9 \anchor{Frenkel p.73}、\anchor{Frenkel p.111} & $N=4$ 扭化、$\MH$、A/B-膜、't Hooft 算子与 Bogomolny 方程 \S3--10 \anchor{src p.16}、\anchor{src p.123} \\
对偶的含义 & 几何 Langlands 猜想（Conjecture 1）\anchor{Frenkel p.65}；阿贝尔情形已证（Deligne）\anchor{Frenkel p.49} & 电磁对偶的物理实现：$S$-对偶在 2d 变成镜像对称（M2）\anchor{src p.71}、\anchor{src p.74} \\
对分歧的处理 & 主场：抛物 $\mathrm{Bun}_{G,y}$、仿射 Hecke 代数、范畴化 \S9.7--9.8 \anchor{Frenkel p.116}、\anchor{Frenkel p.118} & 明确只做非分歧（自陈第一个 gap）\anchor{src p.5} \\
$D$-模的来源 & 共形块 $\to$ 扭曲 $D$-模 \S7.4--7.6 \anchor{Frenkel p.80}；Hecke 特征层作为 Hitchin 量子化的方程 \S9.5 \anchor{Frenkel p.111} & c.c.\ A-膜 $\to$ $D$-模 \S11.1--11.2 \anchor{src p.186}、\anchor{src p.196} \\
群的语言 & 复约化群 $G$ 与 $\Ldual$（起手就是复群）\anchor{Frenkel p.3} & 紧 $G$，复化 $G_\CC$；作者明说这与其他表述不同 \anchor{src p.6} \\
2d CFT 的角色 & 主工具（Part III 全部）\anchor{Frenkel p.73} & 自陈第二个 gap：不理解其用处 \anchor{src p.5} \\
已证范围 & 阿贝尔情形完整；opers 情形由 Beilinson--Drinfeld 等给出 \anchor{Frenkel p.49}、\anchor{Frenkel p.111} & 无定理；物理推导 \\
曲线与边界 & 紧黎曼面 $X$，允许分歧点 \anchor{Frenkel p.116} & 紧 $C$，$g\ge 2$，非分歧 \anchor{src p.5} \\
\bottomrule
\end{tabular}
\end{center}

\subsubsection*{可互译的三处}

两套语言并非互不相干。下面三条是我们能找到的具体接口，按可靠性从高到低排列。

（1）\textbf{电磁交换与自由玻色子 $T$-对偶。} Frenkel 在 \S8.6 用半径为 $\nu$ 与 $1/\nu$ 的自由玻色子 $T$-对偶解释为什么临界 level 仿射代数的中心同构于 $\Ldual$ 的经典 $W$-代数：在 $T$-对偶下 electric 与 magnetic 顶点算子互换（\anchor{Frenkel p.98}）。\pA K--W 的 M2 在数学上正是 Hitchin 纤维化上的 $T$-对偶（SYZ 镜像对称），$\MH(G)$ 与 $\MH(\Ldual)$ 的纤维互为对偶环面（\anchor{src p.71}）。\pA 把这两处 $T$-对偶认作同一机制，是我们的重建：两边各自明写 $T$-对偶，但都没有把这两人写在同一页上。\pB

（2）\textbf{Hecke 特征层与 Hitchin 系统的量子化。} Frenkel \S9.5 证明 $\mathrm{Fun}\,\mathrm{Op}_{\Ldual\mathfrak{g}}(X)\cong D'_{-h^\vee}$，即 Hecke 特征层可以读成 Hitchin 系统量子化后的整体微分方程（\anchor{Frenkel p.111}）。\pA K--W \S11.1 用 $\MH$ 上的 c.c.\ A-膜给出 $\MC$ 上的 $D$-模（\anchor{src p.186}），而 \S11.3 的标题正是「Generalizations of the c.c.\ Brane and Twisted Differential Operators」（\anchor{src p.201}）——说明他们知道这条对译的数学一侧。\pA 于是可以写：\textbf{c.c.\ 膜给出的 $D$-模，就是 Frenkel 书中 Hitchin 量子化给出的那些 $D$-模}。\pB 这条对译的价值在于，它把 2006 最技术的一步（S8）翻译成了 Beilinson--Drinfeld 传统的定理，从而让「2006 没有证明」这句话有了精确含义：缺的不是论证，而是把膜范畴的物理对象等同于 $D$-模层的那一步。\pC

（3）\textbf{分歧：抛物结构与 2-流形算子。} Frenkel 把 tame 分歧处理成 $\mathrm{Bun}_{G,y}$ 上的抛物结构加仿射 Hecke 代数作用，并在 \S9.8 给出局部化函子 $\Delta_y$（\anchor{Frenkel p.118}）。\pA 2006 明确不做分歧（\anchor{src p.5}），但 0905 用 $d-2=2$ 的 2-流形算子加量子参数给出了分歧的 QFT 框架（\anchor{0905.2720 p.6}），并引 Gukov--Witten 2006。\pA 把 Frenkel 的抛物结构与 0905 的 2-流形算子认作同一件事的两侧，是我们的判断：两边都没有写下这个对应，但它解释了为什么「分歧」在两条路线里都表现为「在曲线上加一个点与一个参数」。\pC

\subsubsection*{不能互译的地方：opers}

必须指出一处我们\textbf{找不到}对译的地方。Frenkel 的核心参数化是 $\Ldual$-oper：把 Hecke 特征层的特征值写成 $X$ 上的射影联络，并证明 Hecke 特征层就是该 oper 对应的 $D$-模（\S9.1--9.3，\anchor{Frenkel p.100}）。\pA 2006 里没有 oper 这个对象；0802.0999 的量子参数 $\eta$（\anchor{0802.0999 p.7}）与 opers 的关系是间接的（经 A-膜 $\to$ $D$-模），作者也没有说 ``oper'' 一词。\pA 因此「K--W 侧哪个对象等于 oper」这个问题，在本附录所及的文献里没有明写答案；最接近的猜测是 c.c.\ 膜上的某类扭曲 $D$-模（\anchor{src p.201}），但这仍是猜测。\pC

\subsubsection*{分歧处理是两条路线的分水岭}

K--W 的自陈是：\pA
\begin{origquote}
One obvious gap in our analysis is that we consider only the unramified case of the geometric Langlands program. \anchor{src p.5}
\end{origquote}
这不是疏忽，而是方法论差异。扭化规范论给出的是一套\textbf{非分歧}的算子代数：'t Hooft 算子在光滑点、在平凡丛附近有定义，Hecke 算子也只在非分歧点有干净的几何解释（\anchor{src p.123}）。\pA 分歧需要局部数据（Iwahori 子群、仿射 Hecke 代数、抛物结构），这些正是 Frenkel 的 CFT 机器最擅长的（\anchor{Frenkel p.117}）。\pA 所以两篇的分工不是「谁更完整」，而是「谁处理哪一类点」。\pC 2006 之后，分歧由 Gukov--Witten 2006、Witten 2007（\emph{Gauge Theory and Wild Ramification}）与 0802.0999 的 \S4--5 逐步接手（\anchor{0802.0999 p.5}）——而这条后续路线仍然在扭化规范论的框架内，不是 Frenkel 的 CFT 路线。\pA

\paragraph{自检问题}
\begin{enumerate}
\item Frenkel 的哪一步在 2006 里有对应物，哪一步没有？请分别举一个例子并标出锚点。\pC
\item 为什么「Hitchin 系统量子化」可以作为两条路线的交汇点，而「opers」不能？\pC
\end{enumerate}

% ----------------------------------------------------------------------
\subsection{与 Setter 2013 对照：把 K--W 往下压到 3d 与范畴}
\label{sec:compare-setter}

\subsubsection*{先纠正一个预期}

\begin{kcard}{Setter 的论文实际做了什么}
摘要（\anchor{Setter p.v}）与 \S1.2（\anchor{Setter p.5}）、\S1.3（\anchor{Setter p.7}）一致地表明：这篇论文是\textbf{沿 K--W 的方向推广}，而不是把 K--W 严格化。它把 GL-扭化理论在圆上紧化到 3d，分类 $t=\ii$ 时的面算子与线算子，把 Vafa--Witten 理论实现为「gauged 4d A-模型」的特例，并在附录里分类 2d sigma 模型的各种拓扑缺陷。
\end{kcard}

三条书目事实支撑这个判断。第一，全书 ``Hecke'' 只出现 5 次：4 次在引论中转述 K--W 的结论（例如 \anchor{Setter p.2}），1 次是无关的 Hecke group（Fuchsian 群的一个离散子群）。\pA 第二，全书没有 Hecke 函子、Hecke 作用的定义，也没有把 GL 对应表述成待证命题。\pA 第三，导师是 A. Kapustin，致谢里写「whose ideas are central to this thesis」（\anchor{Setter p.iv}）；这篇论文属于 K--W 学派内部的推进，而不是外部数学家的严格化。\pA 因此本小节的问题要改写为：\textbf{它在哪个层级上做了严格化}，以及它引用了 K--W 的哪些部分作为输入。

\subsubsection*{对照表：输入、推广、严格化、未做}

\begin{center}
\small\textbf{表 A.4}\quad Setter 2013 与 K--W 2006 的分工\\[3pt]
\footnotesize
\begin{tabular}{@{}p{0.13\textwidth}p{0.21\textwidth}p{0.30\textwidth}p{0.13\textwidth}@{}}
\toprule
\tabhead{主题} & \tabhead{K--W 2006} & \tabhead{Setter 2013} & \tabhead{关系} \\
\midrule
4d 扭化与 $t$ & \S3 \anchor{src p.16} & 第 2 章复述扭化与 $t$ 参数（\anchor{Setter p.13}） & 输入 \\
在 $C$ 上紧化 & \S4--5 \anchor{src p.38}、\anchor{src p.56} & 第 8 章把 Vafa--Witten 理论在 $C$ 上紧化到 $\MH$ 上的 A-模型（\anchor{Setter p.35}） & 推广 \\
膜范畴 & \S5.6 \anchor{src p.77}（隐式使用） & \S4.1--4.2 显式给出 2d TFT 的范畴与 2-范畴（\anchor{Setter p.33}） & 严格化（一般 TFT） \\
线算子 & \S6 \anchor{src p.78} & 第 6 章确定 $t=\ii$ 的完整幺半线算子范畴，含 Wilson 线为特例（\anchor{Setter p.7}） & 精确化 $+$ 推广 \\
面算子 & 未系统处理 & 第 6 章分类 $t=\ii$ 的面算子 2-范畴（module category 描述），并处理 $t=0,1$ 的阿贝尔情形 & 新增 \\
't Hooft $=$ Hecke & \S9--10 \anchor{src p.123}、\anchor{src p.148} & 不处理 & 未做 \\
c.c.\ 膜$\to$\mbox{$D$-模} & \S11 \anchor{src p.184} & 不处理 & 未做 \\
$n$-范畴层级 & 隐式（膜范畴） & \S4.3--4.4：面算子 2-范畴 $=$ 3d TFT 的边界条件 2-范畴（\anchor{Setter p.39}、\anchor{Setter p.41}） & 严格化 \\
Vafa--Witten & 未涉及 & 第 3 章：作为 gauged 4d A-模型的特例（\anchor{Setter p.20}） & 新增 \\
定理化 & 无 & 22 个 Proposition，集中在第 9 章与附录（2d sigma 模型缺陷、topological bibrane） & 局部定理化 \\
\bottomrule
\end{tabular}
\end{center}

\subsubsection*{严格化的真实层级：范畴，不是定理}

Setter 的严格化确实存在，但发生在\textbf{范畴层}而非定理层。\S4.1 开头把范畴定义为「多对象的代数」（\anchor{Setter p.33}）；\S4.4 用维数约化把 4d 中的面算子 2-范畴实现为 3d TFT 的边界条件 2-范畴（\anchor{Setter p.41}）；作者自己把这条路线概括成「(1.1) 是集合的等价，(1.2) 是范畴的等价，(1.3) 是 2-范畴的等价」，并说每升一级就「加入额外的对象、显式化此前隐含的刚性」（\anchor{Setter p.5}）。\pA 他自己也划了界：\pA
\begin{origquote}
we stop at the 2-categorical level in this thesis and leave the braided, monoidal structure to future work. \anchor{Setter p.5}
\end{origquote}
这与「把 K--W 的物理论证放进 Atiyah 型 TFT 框架」不是一回事：Setter 组织材料时确实用「TFT 给范畴」的语言，但全书没有引用 Atiyah 的 TFT 公理，也没有把 GL 表述成从协边范畴出发的函子。\pC 他严格化的是\textbf{缺陷的范畴结构}——面算子、线算子、边界条件的层级——而不是 GL 对应本身。\pC

\subsubsection*{依赖顺序：整篇为输入，不重证任何一步}

文献表里 K--W 只有一条整条引用 [1]，正文 38 处引用都不带节号。\pA 依赖顺序是：第 2 章复述 2006 的扭化与 $t$（\anchor{Setter p.13}）；\S1.1 复述 2006 的 $\Sigma\times C$ 图像与「收缩 $C$」极限（\anchor{Setter p.3}）；第 5 章从 2006 的 4d TFT 出发再做圆紧化；结果列表明说 $t=\ii$ 的线算子范畴包含 [1] 的 Wilson 线为特例（\anchor{Setter p.7}）。\pA 换句话说，Setter 把 K--W 的 S2--S5 整体当作已经成立的物理结论，自己从那里\textbf{往下}走到 3d 与 2d。\pB 这也回答了「哪些地方补了证明」：他补的是 3d 有效理论的构造、面算子 2-范畴的 module category 描述、线算子幺半范畴的确定——这些在 K--W 里没有；他没有重证 K--W 的任何一步。\pB

\subsubsection*{与 2009 Witten 的方向对比}

两篇都以 2006 为输入，但方向相反：Witten 2009 把起点\textbf{往上抬}到六维，用 $(0,2)$ 理论解释 4d 对偶为什么存在；Setter 把同一套理论\textbf{往下压}到三维与二维，用 2-范畴语言整理它的缺陷。\pC 两者的共同点是都没有回到 2006 的八步链去补缺口：2009 不重做 S2--S8，Setter 不重做 S9--S11。\pC 因此对本书读者，二者的用途也不同：想理解「为什么」，读 2009 的 §3--4；想理解「缺陷的范畴结构」，读 Setter 的第 4、6 章；想理解 2006 的主链，两者都不能替代底本。\pC 至于「K--W 推导的数学严格化」，它不在 Setter 这里——真正把 2006 的某些步骤变成定理的，是 Beilinson--Drinfeld--Frenkel 传统（Frenkel \S9.4--9.5 的 Hecke 性质与 Hitchin 量子化），以及 0802.0999 提供的 Hitchin 几何表述。\pC

\paragraph{自检问题}
\begin{enumerate}
\item Setter 的论文里，哪一章最接近「把 K--W 严格化」？它的严格化对象是什么？\pC
\item Setter 与 Witten 2009 都以 2006 为输入，为什么不能把二者互相替代？\pC
\end{enumerate}

% ----------------------------------------------------------------------
\subsection{小结：三份文献各占主链哪一段}
\label{sec:compare-summary}

\begin{center}
\small\textbf{表 A.5}\quad 三份文献的位置与替代边界\\[3pt]
\footnotesize
\begin{tabular}{@{}p{0.14\textwidth}p{0.38\textwidth}p{0.38\textwidth}@{}}
\toprule
\tabhead{文献} & \tabhead{在 2006 主链上的位置} & \tabhead{它不能替代 2006 的什么} \\
\midrule
Witten 2008 & S3--S4 的几何化：Hitchin 纤维化、镜像对称、分歧与野分歧（\S4--5）\anchor{0802.0999 p.5}、\anchor{0802.0999 p.7} & 不能替代 4d 规范论的推导；GL 的物理想法仍引 2006 \anchor{0802.0999 p.14} \\
Witten 2009 & 换起点到六维：S1 与 S7 各压成一句，新增 $(0,2)$、M5 膜、复曲面版 GL & 不能替代 S2--S8 的技术内容；未补齐分歧、非 simply-laced 与证明 \\
Frenkel 2003 & 换入口到 CFT：S8 的数学对应物（$D$-模与 Hitchin 量子化）、分歧的主场 & 不能替代 2006 的对偶来源；不处理 4d 规范论与膜 \\
Setter 2013 & 换维度到 3d/2d：S3--S5 的范畴化、面算子与线算子的完整分类 & 不能替代 2006 的 S7--S11；不定义 Hecke 作用、不给 GL 定理 \\
\bottomrule
\end{tabular}
\end{center}

把三份文献叠在 2006 的主链上，可以看到一个整齐的分工：2008 补 S3--S4 的几何，2009 换 S1 与 S7 的起点，Frenkel 补 S8 的数学对应，Setter 补 S5 之后的范畴结构。\pC 四者都没有做的，是 2006 的第二 gap（2d CFT 与规范论的统一）、非阿贝尔 GL 的证明、非 simply-laced 的完整处理，以及野分歧的范畴化。\pC 这恰好说明 2006 那篇 231 页的价值不在任何单一步骤的严格性，而在它把一条从 4d 规范论到 Hecke 特征层的\textbf{完整路线}第一次画了出来；后续文献各取一段深化，却没有人重画整条路线。\pC 本书正文按 S1--S8 重排这条路线，附录 B 提供配套语料，附录 C 提供讨论课问题集，正是为了让读者能同时看到「路线」与「各段的后续」。\pC

\begin{keyidea}
一句话总结：2008 把 2006 的几何讲清楚，2009 把 2006 的起点换掉，Frenkel 把 2006 的终点用 CFT 重造，Setter 把 2006 的框架往下压成范畴语言。四者都是 2006 的\textbf{后继}，没有一个是 2006 的\textbf{替代}。
\end{keyidea}
